\chapter{\texorpdfstring{Rational parameters $a = r/s$}{Rational parameters a = r/s}}
\label{chap:rational}

\emph{This chapter is a plan, not a record.} Nothing in it is formalized yet. It states the
rational version of the Section~3 theorem obtained in the working notes
(\texttt{tex/stoll-method-extensions.tex}: Theorems R${}^+$ and R${}^-$ there), reduces it to
the material of Chapters~1--5 plus the new lemmas listed here, and fixes the order of work.
Every node below is uncoloured; the arrows into it are the point. Proposed declaration names
are given in \lname{typewriter}; they follow the naming grammar of the project.

Throughout, $a = r/s \in \QQ$ with $r, s \in \ZZ$ coprime, $s \ge 1$, $r \ne 0$, and $-rs$ not a
square (else $f = X^2 + a$ is reducible). The iterates $f_n$, the fields $K_n$ and the groups
$\Om_n$ are defined as in Chapter~0 with $a \in \QQ$.

\section{The Galois side for a rational parameter}

The formalized Section~1 (Chapter~1) is parametrized by an integer $a$ and uses the standing
hypothesis ``$-a$ is not a square'' in exactly three ways: to know that every $f_n$ is
irreducible (Corollary~\ref{cor:1.3}, a size argument that is false for rational $a$: for
$a = 1/3$, $c_2 = 4/9$ is a square while $f_2$ is irreducible), to know $c_{n+1} \ne 0$, and
through the integer factors $b_n$ of Lemma~\ref{lem:1.1}. Only the first two matter for Lemmas
\ref{lem:1.4}--\ref{lem:1.6} and the equivalence (a)$\iff$(b) of Theorem~\ref{thm:section1}.

\begin{theorem}[Section 1 for rational $a$]
  \label{thm:rat-section1}
  \uses{def:galois, def:wreath, def:twoindep}
  \lean{QuadraticIterates.section1_tfae_of_irreducible,
    QuadraticIterates.section1_a_iff_b_of_irreducible, QuadraticIterates.section1_b_iff_c}
  \leanok
  Let $a \in \QQ^\times$ and suppose that $f_k$ is irreducible over $\QQ$ for all $k \le n$.
  Then $\Om_n \cong \Ct{n}$ if and only if $c_1, \dots, c_n$ are $2$-independent in
  $\QQ^\times/(\QQ^\times)^2$. In particular (Lemma~\ref{lem:1.2}), if none of $c_1, \dots, c_n$ is
  a square, then $\Om_n \cong \Ct{n}$ iff $c_1, \dots, c_n$ are $2$-independent.
\end{theorem}
\begin{proof}
  \uses{lem:1.4, lem:1.5, lem:1.6, lem:1.2}
  \leanok
  The induction of Theorem~\ref{thm:section1}, (a)$\iff$(b), verbatim: Lemma~\ref{lem:1.4}
  reduces level $n+1$ to level $n$ and the degree $[K_{n+1} : K_n] = 2^{2^n}$; Lemma~\ref{lem:1.6}
  turns the degree into ``$c_{n+1}$ is not a square in $K_n$'' (its proof uses only that $f_n$ is
  irreducible, hence separable with $2^n$ roots, and $c_{n+1} = f_n(a) \ne 0$, which follows from
  the irreducibility of $f_{n+1}$: otherwise $X \mid f_{n+1}$); Lemma~\ref{lem:1.5} turns that
  into $2$-independence. If no $c_k$ is a square, all $f_k$ are irreducible by Lemma~\ref{lem:1.2}.
  \formalnote{Done as planned (2026-09-07, with the author's approval of the change of the
  comparator challenge): \lname{splittingField}, \lname{GaloisGroup}, \lname{rootShift} take a
  rational $a$; \lname{cSeq} is polymorphic in a commutative ring and \lname{bSeq} in a domain,
  $c$ being the $\gamma$-sequence of $X^2 + a$ with $\varepsilon = -1$ over that ring, with the
  cast lemmas \lname{intCast\_cSeq}, \lname{ratCast\_cSeq}, \lname{intCast\_bSeq} and
  \lname{bSeq\_eq\_moebiusFactorK} (over a field, $b_n$ is the Möbius product itself). Since
  the expected type propagates through the polymorphic definitions, the integer statements
  \lname{section1\_a\_iff\_b}, \lname{section1\_tfae},
  \lname{section1\_of\_not\_isSquare\_abs\_bSeq} keep their text and are the rational
  statements at $a \in \ZZ \subseteq \QQ$, proved via Corollary~\ref{cor:1.3}; the same holds
  for the theorems of Chapters~3 and~5 and for the comparator challenges. Lemma~\ref{lem:1.6}
  is stated with $f_n$ irreducible and $c_{n+1} \ne 0$
  (\lname{cSeq\_succ\_ne\_zero\_of\_irreducible} from $f_{n+1}$ irreducible);
  Lemma~\ref{lem:1.5} needs no hypothesis on $a$ at all, because $c_{k+1}$ is a square in
  $K_n$ for every $a$ (the norm identity with multiplicity). ``b) $\Leftrightarrow$ c)''
  (\lname{section1\_b\_iff\_c}) needs only that no $c_k$ vanishes. The growth argument of
  \lname{Irreducibility.lean} stays integral.}
\end{proof}

\section{The rescaled sequence with denominators}

\begin{definition}[the $w$-sequence]
  \label{def:wseq}
  \uses{def:cseq}
  \lean{QuadraticIterates.wSeq, QuadraticIterates.wSeq_succ, QuadraticIterates.intCast_wSeq,
    QuadraticIterates.wSeq_pos_of_pos, QuadraticIterates.wSeq_pos_of_neg_lt,
    QuadraticIterates.pow_le_wSeq_of_two_mul_le, QuadraticIterates.wSeq_pos_of_two_mul_le}
  \leanok
  For coprime $r, s \in \ZZ$, $s \ge 1$, and a sign $\varepsilon = \pm 1$, let $w_1 = 1$ and
  $w_{n+1} = r w_n^2 + \varepsilon s^{2^n - 1}$ (junk value $w_0 = 0$). For the $\gamma$-sequence
  of $aX^2 + \varepsilon$ with sign $\varepsilon$, i.e.\ $\gamma_1 = 1$,
  $\gamma_{n+1} = a\gamma_n^2 + \varepsilon$, one has $\gamma_n = w_n / s^{2^{n-1} - 1}$ in every
  commutative ring in which $s$ is invertible. For $\varepsilon = 1$ this is $\gamma_n = f^n(0)/a$
  and $c_n = a \gamma_n = r w_n / s^{2^{n-1}}$ for $n \ge 1$; $\varepsilon = -1$ with $|r|$ in
  place of $r$ is Stoll's positive normalization for $a \le -2$.
  \formalnote{\lname{wSeq r s ε}, a function of $n$ with values in $\ZZ$, by recursion
  (\lname{wSeq\_succ}); the identity is \lname{intCast\_wSeq}: in a commutative ring $S$ with an
  element $u$, $s u = 1$, the cast of $w_n$ is $s^{2^{n-1}-1}\gamma_n$ for the $\gamma$-sequence of
  $(r u) X^2 + \varepsilon$ with sign $\varepsilon$, an induction of a dozen lines, used over
  $\ZZ[1/s]$ (Lemma~\ref{lem:gamma-away}), over $\QQ$ (Lemma~\ref{lem:rat-sqclass}) and over
  $\ZZ/M\ZZ$ (Lemma~\ref{lem:w-gamma-zmod}). Because of the varying term $s^{2^n-1}$, $w$ is not
  itself a $\gamma$-sequence over $\ZZ$; it is one over $\ZZ[1/s]$ up to units, which is how the
  general theory of Chapter~2 enters.}
\end{definition}

\begin{lemma}[integrality and coprimality]
  \label{lem:w-coprime}
  \uses{def:wseq}
  \lean{QuadraticIterates.isCoprime_wSeq_left, QuadraticIterates.isCoprime_wSeq_right}
  \leanok
  $w_n \in \ZZ$ and $\gcd(w_n, rs) = 1$ for all $n \ge 1$.
\end{lemma}
\begin{proof}
  \leanok
  $w_{n+1} \equiv r w_n^2 \pmod s$ and $w_{n+1} \equiv \varepsilon s^{2^n - 1} \pmod r$; induction
  from $w_1 = 1$ and $\gcd(r, s) = 1$.
  \formalnote{\lname{isCoprime\_wSeq\_left} ($w_n$ coprime to $r$, no induction needed) and
  \lname{isCoprime\_wSeq\_right} ($w_n$ coprime to $s$, by \lname{Nat.le\_induction}), both from
  \lname{IsCoprime.add\_mul\_left\_left}.}
\end{proof}

\begin{lemma}[{$w$ is a $\gamma$-sequence over $\ZZ[1/s]$, up to units}]
  \label{lem:gamma-away}
  \uses{def:wseq, def:gammaseq, lem:2.1, thm:valshape, lem:moebiusfactor}
  \lean{QuadraticIterates.aAway, QuadraticIterates.normPolyAway, QuadraticIterates.wSeqAway,
    QuadraticIterates.normPolyAway_eq, QuadraticIterates.eval_zero_normPolyAway,
    QuadraticIterates.gammaSeq_normPolyAway_one, QuadraticIterates.wSeqAway_eq_mul_gammaSeq,
    QuadraticIterates.wSeqAway_two, QuadraticIterates.isRelPrime_moebiusFactorR_wSeqAway}
  \leanok
  Let $R = \ZZ[1/s]$, a unique factorization domain whose primes are the primes $p \nmid s$, and
  let $a = r/s \in R$. The $\gamma$-sequence of $aX^2 + 1$ over $R$ satisfies
  $w_n = s^{2^{n-1}-1}\gamma_n$ in $R$, with $s$ a unit of $R$. Hence $(w_n)$, viewed in $R$, is
  a strong divisibility sequence (Lemma~\ref{lem:2.1}) with the constant-valuation shape
  (Theorem~\ref{thm:valshape}) at every prime $p \nmid s$: $p \mid w_n$ iff $d_0(p) \mid n$, where
  $d_0(p)$ is the least index with $p \mid w_{d_0}$, and $v_p(w_{kd_0}) = v_p(w_{d_0})$ for all
  $k \ge 1$. Consequently the Möbius factors $\prod_{d \mid n} w_d^{\mu(n/d)}$ lie in $R$ and are
  pairwise relatively prime in $R$ (Lemma~\ref{lem:moebiusfactor}, Theorem~\ref{thm:valshape}).
\end{lemma}
\begin{proof}
  \uses{lem:2.1, thm:valshape, lem:moebiusfactor}
  \leanok
  The identity is the case $S = R$ of the identity in Definition~\ref{def:wseq}. Multiplying a
  sequence termwise by units changes neither $\gcd$'s up to association nor valuations, so
  Lemma~\ref{lem:2.1} and Theorem~\ref{thm:valshape} for the $\gamma$-sequence transfer to
  $(w_n)$, and Lemma~\ref{lem:moebiusfactor} and Theorem~\ref{thm:valshape} give the Möbius
  factors.
  \formalnote{Since Mathlib v4.34.0-rc2 a localization of a unique factorization monoid is one
  (\lname{UniqueFactorizationMonoid (Localization S)}); \lname{Localization.Away s} is a domain
  since $s \ne 0$ (\lname{Localization.Away.isDomain}, to be registered as an instance for
  \lname{[NeZero s]}). The API of Chapter~2 (\lname{gammaSeq\_associated\_gcd},
  \lname{factorization\_gammaSeq\_shape}) and \lname{moebiusFactorR\_isRelPrime} need a
  \lname{NormalizedGCDMonoid} instance, a \lname{DecidableEq} instance and normalized primes;
  Mathlib has neither instance for the localization (its \lname{DecidableEq} instance needs
  cancellation on the base monoid, which fails for the ring $\ZZ$). Both are supplied
  \emph{locally}: \lname{UniqueFactorizationMonoid.strongNormalizationMonoid} and
  \lname{toNormalizedGCDMonoid} as local instances of the one file that applies the API (an
  arbitrary normalization, visible in no exported statement, since relative primality and
  divisibility do not mention it), and \lname{classical} inside the proofs. The inverse of $s$ is
  \lname{IsLocalization.Away.invSelf}. Declarations: \lname{aAway} ($a = r/s$ in $R$),
  \lname{normPolyAway} ($aX^2 + \varepsilon$ over $R$), \lname{wSeqAway} (the sequence
  $n \mapsto w_n$ in $R$), \lname{wSeqAway\_eq\_mul\_gammaSeq}, the private
  \lname{wSeqAway\_associated\_gcd} and \lname{factorization\_wSeqAway\_shape}, and
  \lname{isRelPrime\_moebiusFactorR\_wSeqAway}, each a few lines, exactly as
  \lname{factorization\_cSeq\_shape} and \lname{isCoprime\_bSeq} in Chapter~1.}
\end{proof}

\paragraph{The alternative not taken.}
One can avoid $\ZZ[1/s]$ and prove the shape of $w$ over $\ZZ$ directly: the twisted rigidity
$w_d^2 \mid w_{d+j} - s^{2^{j-1}(2^d-1)} w_j$ (an induction like that of Lemma~\ref{lem:2.1}),
then the Euclidean and periodic-shape arguments of Chapter~4 with the twist $s^{(\cdot)}$, a
unit modulo any $p \nmid s$; this needs twisted generalizations of the two supporting lemmas
(\lname{gcd\_eq\_normalize\_of\_dvd\_sub}, \lname{factorization\_periodic\_shape}) with a unit
factor in the congruence, about a hundred lines against about a hundred and fifty for the
localization route, whose surplus consists of general lemmas on localizations (a domain instance,
the transfer of primes and of divisibility) and of the descent below. The localization route is
chosen because it reuses Chapter~2 verbatim instead of re-deriving it with a twist, and because
its new lemmas are of Mathlib generality.

\begin{definition}[the integer Möbius factors]
  \label{def:betaint}
  \uses{def:wseq}
  \lean{QuadraticIterates.betaInt, QuadraticIterates.betaInt_eq_moebiusFactorR}
  \leanok
  $\beta_n := \prod_{d \mid n} w_d^{\mu(n/d)}$ and $E(n) := \sum_{d \mid n} (2^{d-1} - 1)\mu(n/d)$,
  so that $\beta_n = s^{E(n)} \prod_{d \mid n} \gamma_d^{\mu(n/d)}$ and $E(n) \equiv \mu(n) \pmod 2$
  for $n \ge 2$.
  \formalnote{\lname{betaInt r s ε n := moebiusFactorR (wSeq r s ε) n}, the same device as
  \lname{bSeq}: \lname{moebiusFactorR} is defined for any sequence in a domain, and its API
  (\lname{moebiusFactorR\_mul\_denProd}, \lname{algebraMap\_moebiusFactorR},
  \lname{prod\_moebiusFactorR}, \lname{moebiusFactorR\_ne\_zero}) uses strong divisibility only
  through $D \mid N$ (Lemma~\ref{lem:moebiusfactor}); these four lemmas have primary forms
  \lname{\_of\_dvd} with the hypothesis $D \mid N$, of which the strong-divisibility forms are
  corollaries, and \lname{map\_moebiusFactorR}: for $D \mid N$ and a ring homomorphism $\varphi$
  into a domain killing no $w_d$, $\varphi(\beta_n)$ is the Möbius factor of $\varphi \circ w$,
  both being the solution of $\beta D = N$. Still to come: \lname{twistExp}, the integer $E(n)$
  as a sum over \lname{n.divisorsAntidiagonal}, with \lname{twistExp\_emod\_two}
  ($E(n) \equiv \mu(n) \bmod 2$ for $n \ge 2$).}
\end{definition}

\begin{lemma}[the $\beta_n$ are integers, pairwise coprime, coprime to $rs$]
  \label{lem:betaint}
  \uses{def:betaint, lem:gamma-away, lem:w-coprime}
  \lean{QuadraticIterates.denProd_wSeq_dvd_numProd_wSeq, QuadraticIterates.betaInt_mul_denProd,
    QuadraticIterates.intCast_betaInt, QuadraticIterates.betaInt_ne_zero,
    QuadraticIterates.wSeq_eq_prod_betaInt, QuadraticIterates.isCoprime_betaInt,
    QuadraticIterates.isCoprime_betaInt_left, QuadraticIterates.isCoprime_betaInt_right,
    QuadraticIterates.betaInt_pos, IsLocalization.Away.isRelPrime_of_isRelPrime_algebraMap,
    IsLocalization.Away.dvd_of_algebraMap_dvd_of_isCoprime,
    IsLocalization.Away.prime_algebraMap_of_not_dvd, IsLocalization.prime_algebraMap_of_prime,
    IsLocalization.dvd_of_algebraMap_dvd}
  \leanok
  For $n \ge 2$, $\beta_n \in \ZZ$, $\gcd(\beta_n, rs) = 1$, $\gcd(\beta_m, \beta_n) = 1$ for
  $m \ne n$, and $v_p(\beta_n) = v_p(w_{d_0(p)})\,[n = d_0(p)]$ for every prime $p \nmid rs$.
\end{lemma}
\begin{proof}
  \uses{lem:gamma-away, lem:w-coprime}
  \leanok
  Everything holds in $R = \ZZ[1/s]$ by Lemma~\ref{lem:gamma-away} and descends to $\ZZ$: with
  $N = \prod_{\mu(n/d) = 1} w_d$ and $D = \prod_{\mu(n/d) = -1} w_d$, both coprime to $rs$
  (Lemma~\ref{lem:w-coprime}), $D \mid N$ in $R$ means $Dm = Ns^k$ in $\ZZ$ for some $m$, $k$,
  so $D \mid N$ in $\ZZ$ and $\beta_n = N/D \in \ZZ$ is coprime to $rs$. A prime $p$ dividing
  $\beta_m$ and $\beta_n$ is prime to $s$, hence a prime of $R$ dividing both Möbius factors
  there, which are relatively prime. The valuation formula is that of Lemma~\ref{lem:gamma-away}.
  \formalnote{The descent rests on Mathlib's transport lemmas for localization maps
  (\lname{Mathlib/GroupTheory/MonoidLocalization/Divisibility.lean} and
  \lname{UniqueFactorization.lean}, from the same PR \#33832 as the instance):
  \lname{Submonoid.LocalizationMap.map\_dvd\_map} ($\varphi(a) \mid \varphi(b)$ iff $a \mid mb$ for
  some $m$ in the submonoid), \lname{map\_isUnit\_iff} and \lname{map\_prime}. Three lemmas of
  Mathlib generality are added on top of them, in the mirror file of
  \lname{Mathlib/RingTheory/Localization/Away/Basic.lean}: \lname{IsLocalization.Away.dvd\_of\_algebraMap\_dvd\_of\_isCoprime}
  ($\varphi(D) \mid \varphi(N)$ and $D$ coprime to $x$ give $D \mid N$, three lines from
  \lname{map\_dvd\_map}), \lname{IsLocalization.prime\_algebraMap\_of\_prime} --- stated
  \emph{exactly} as in the open Mathlib PR \#37247 (Yuma Mizuno, approved in April 2026,
  awaiting a rebase since July), which adds this ring-level wrapper of \lname{map\_prime} together
  with the prime-divisibility transport \lname{algebraMap\_dvd\_algebraMap\_iff\_of\_prime}; the
  local copy is deleted when that PR lands --- and its consequence
  \lname{IsLocalization.Away.prime\_algebraMap\_of\_not\_dvd} (a prime $p \nmid x$ stays prime,
  via \lname{IsLocalization.Away.algebraMap\_isUnit\_iff}). No PR adds the domain instance for
  \lname{Localization.Away x} or a normalization on localizations. Finally
  \lname{IsLocalization.Away.isRelPrime\_of\_isRelPrime\_algebraMap}: in a UFD, elements the
  first of which is coprime to $x$ and whose images are relatively prime are relatively prime,
  since a common prime factor would stay prime. Then \lname{denProd\_wSeq\_dvd\_numProd\_wSeq},
  and the API of \lname{betaInt} (\lname{betaInt\_mul\_denProd}, \lname{intCast\_betaInt},
  \lname{betaInt\_ne\_zero}, \lname{wSeq\_eq\_prod\_betaInt}) from the \lname{\_of\_dvd} forms;
  \lname{isCoprime\_betaInt} from the transfer lemma, \lname{map\_moebiusFactorR} and
  \lname{isRelPrime\_moebiusFactorR\_wSeqAway}; \lname{isCoprime\_betaInt\_left} and
  \lname{isCoprime\_betaInt\_right} ($\beta_n$ coprime to $r$ and to $s$) from $\beta_n D = N$. The
  hypotheses throughout are $s \ne 0$, $\gcd(r, s) = 1$, $\varepsilon^2 = 1$ and $w_n \ne 0$ for
  $n \ge 1$, the last one supplied later by positivity.}
\end{proof}

\begin{lemma}[square classes]
  \label{lem:rat-sqclass}
  \uses{def:betaint, def:bseq}
  In $\QQ^\times/(\QQ^\times)^2$: $c_1 = b_1 \sim -rs$ and $b_n \sim (-s)^{\mu(n)} \beta_n$ for
  $n \ge 2$.
\end{lemma}
\begin{proof}
  $c_d = r w_d / s^{2^{d-1}} \sim r w_d$ for $d \ge 2$ (even exponent), so
  $b_n \sim (-rs)^{\mu(n)} \prod_{d \mid n, d \ge 2} (r w_d)^{\mu(n/d)} = (-s)^{\mu(n)} r^{\sum_{d \mid n}
  \mu(n/d)} \beta_n$ with $w_1 = 1$ and $\sum_{d \mid n} \mu(n/d) = 0$.
  \formalnote{Stated with \lname{sqClass} (\lname{QuadraticIterates.Mathlib.FieldTheory.Multiquadratic}),
  additively: \lname{sqClass\_bSeq\_eq}, $[b_n] = \mu(n)\,[-s] + [\beta_n]$; the proof is
  \lname{sqClass\_bSeq\_eq\_sum\_divisorsAntidiagonal} plus \lname{sqClass\_cSeq\_eq},
  $[c_d] = [r w_d]$ for $d \ge 2$.}
\end{proof}

\begin{lemma}[shared-part lemma]
  \label{lem:rat-shared}
  \uses{lem:rat-sqclass, lem:betaint, thm:rat-section1}
  If $-rs$ is not a square and $|\beta_n|$ is not a square for every $n \ge 2$, then
  $c_1, \dots, c_n$ are $2$-independent for every $n$, hence $\Om_n \cong \Ct{n}$ for all $n$.
\end{lemma}
\begin{proof}
  \uses{lem:betaint, lem:rat-sqclass}
  A nonempty $S$ with $\prod_{n \in S} b_n$ a square has, by pairwise coprimality and coprimality to
  $rs$, $|\beta_n|$ a square for each $n \in S$, $n \ge 2$; so $S = \{1\}$, but $b_1 \sim -rs$ is not
  a square. The $b$'s are $2$-independent, hence the $c$'s (Möbius inversion in $\QQ^\times$), and
  no $c_k$ is a square, so Theorem~\ref{thm:rat-section1} applies.
  \formalnote{\lname{nonempty\_mulEquiv\_of\_forall\_not\_isSquare\_abs\_betaInt}, the rational
  analogue of \lname{section1\_of\_not\_isSquare\_abs\_bSeq}. The integer argument uses
  \lname{twoIndependent\_iff\_of\_pairwise\_isCoprime}; here the family
  $i \mapsto (-s)^{\mu(i+1)} \beta_{i+1}$ is not pairwise coprime, so the lemma to use is: for
  pairwise coprime $\beta_i$ coprime to $u$, if $u \prod_{i \in S} \beta_i$ is a square then each
  $|\beta_i|$, $i \in S$, is a square
  (\lname{Int.isSquare\_abs\_of\_isSquare\_mul\_prod\_of\_pairwise\_isCoprime} in
  \lname{QuadraticIterates/Mathlib/Algebra/Squares.lean}, from
  \lname{Int.isSquare\_abs\_of\_isSquare\_mul\_of\_isCoprime}; here $u = (-s)^e$).}
\end{proof}

\section{Reflection with the twist}
\label{sec:rat-reflection}

\begin{lemma}[$w$ and $\gamma$ modulo $M$]
  \label{lem:w-gamma-zmod}
  \uses{def:wseq, def:gammaseq}
  \lean{QuadraticIterates.intCast_wSeq}
  \leanok
  Let $M \ge 1$ be coprime to $s$, and let $\bar a = \bar r \bar s^{-1} \in \ZZ/M\ZZ$. Then
  $\bar w_n = \bar s^{\,2^{n-1} - 1} \gamma_n(\bar a)$ in $\ZZ/M\ZZ$, where $\gamma(\bar a)$ is the
  $\gamma$-sequence of $\bar a X^2 + 1$ (with $\varepsilon = 1$) over $\ZZ/M\ZZ$.
\end{lemma}
\begin{proof}
  \leanok
  Induction on $n$ from the recursions.
  \formalnote{The instance of \lname{intCast\_wSeq} (Definition~\ref{def:wseq})
  at $S = \ZZ/M\ZZ$, $\varphi$ the cast and $u$ the inverse of the unit $\bar s$; the
  $\gamma$-sequence over \lname{ZMod M} is \lname{gammaSeq} of $\bar a X^2 + 1$ with
  $\bar a = \bar r \bar s^{-1}$.}
\end{proof}

\begin{lemma}[reflection with the twist]
  \label{lem:rat-reflection}
  \uses{lem:w-gamma-zmod, def:betaint, lem:betaint}
  \lean{QuadraticIterates.not_isSquare_betaInt_of_dvd_add_succ,
    QuadraticIterates.not_isSquare_betaInt_of_dvd_add_succ_of_two_le,
    QuadraticIterates.not_isSquare_betaInt_of_squarefree_of_dvd,
    QuadraticIterates.intCast_betaInt_eq_div, QuadraticIterates.reflNum,
    QuadraticIterates.isCoprime_reflNum_right, QuadraticIterates.reflNum_pos,
    QuadraticIterates.reflNum_one}
  \leanok
  Let $n \ge 2$, $n_0 = n/\rad(n)$, and $M$ coprime to $s$ and to $w_{n_0}$. If $M$ divides the
  numerator $w_{n_0} s^{2^{n_0-1}} + w_{n_0+1}$ of $\gamma_{n_0} + \gamma_{n_0+1}$, then
  $\beta_n \equiv -s^{E(n)} u^2 \pmod M$ for a unit $u$. Hence $|\beta_n|$ is not a square if
  $\beta_n > 0$ and $-1$ (for $n$ not squarefree) resp.\ $-s$ (for $n$ squarefree) is a non-square
  modulo $M$.
\end{lemma}
\begin{proof}
  \uses{lem:w-gamma-zmod, lem:betaint, lem:2.1, lem:2.2, lem:moebiussign, lem:zmodsquare}
  \leanok
  In $\ZZ/M\ZZ$, $\gamma_{n_0+1} = -\gamma_{n_0}$, so $\gamma_{n_0 t} = -\gamma_{n_0}$ for all
  $t \ge 2$ (Lemma~\ref{lem:li-succ}: \lname{gammaSeq\_mul\_eq\_two\_mul} and evenness), and
  $\prod_{t \mid \rad n} \gamma_{n_0 t}^{\mu(\rad(n)/t)} = -u^2$ (\lname{prod\_gammaSeq\_mul\_eq\_neg\_prod},
  \lname{isUnit\_prod\_gammaSeq\_mul}). Multiply by $s^{E(n)}$ (Lemma~\ref{lem:w-gamma-zmod}).
  For $n$ not squarefree $E(n)$ is even, for $n$ squarefree odd.
  \formalnote{\lname{not\_isSquare\_betaInt\_of\_dvd\_add\_succ}: hypotheses $n \ge 2$,
  $n' = \rad(n)$, $n = k n'$, an element $u$ of $\ZZ/M\ZZ$ with $s u = 1$ (the inverse of the unit
  $\bar s$, passed as a certificate), $M \mid N_k$ (\lname{reflNum r s ε k}, the numerator
  $w_k s^{2^{k-1}} + w_{k+1}$, coprime to $s$ and positive when $w$ is), and $-\bar s^{F}$ not a
  square in \lname{ZMod M}, where $F = \sum_{t \mid n'} (2^{kt-1} - 1)$ is the exponent the
  denominators produce ($F \equiv E(n) \bmod 2$); conclusion $\beta_n$ is not a square in $\ZZ$.
  The proof: \lname{intCast\_betaInt\_eq\_div} writes $\beta_n = P/Q$ with $P$, $Q$ the products
  of the $w_{kt}$ over the two halves of the divisors of $n'$; modulo $M$ these are
  $\bar s^{F^\pm}$ times the corresponding products of $\gamma_{kt}$
  (\lname{intCast\_wSeq}), which are negatives of each other and units
  (\lname{prod\_gammaSeq\_mul\_eq\_neg\_prod}, \lname{isUnit\_prod\_gammaSeq\_mul} over
  \lname{ZMod M}, from $\gamma_k + \gamma_{k+1} = 0$ and \lname{gammaSeq\_add\_succ\_dvd},
  \lname{isCoprime\_gammaSeq\_add\_succ}); \lname{ZMod.isSquare\_mul\_of\_isSquare\_div} (the
  generalization of \lname{ZMod.isSquare\_neg\_one\_of\_isSquare\_div} to a twist) then makes
  $-\bar s^{F^+ + F^-}$ a square. Coprimality of $M$ with $w_k$ is automatic. The two consumed
  forms: \lname{not\_isSquare\_betaInt\_of\_dvd\_add\_succ\_of\_two\_le} ($k \ge 2$, $F$ even, so
  $-1$ not a square suffices) and \lname{not\_isSquare\_betaInt\_of\_squarefree\_of\_dvd}
  ($k = 1$, $F$ odd, $M \mid r + (1 + \varepsilon) s$, so $-\bar s$ not a square suffices), by
  \lname{IsUnit.isSquare\_mul\_sq\_iff}. Both sign worlds are the same statements with
  $(r, \varepsilon) = (|r|, -1)$; positivity of $\beta_n$ is not needed for
  ``$\beta_n$ is not a square'' but for passing to $|\beta_n|$.}
\end{proof}

\section{The residue classes}

\begin{proposition}[the classes (C${}^\pm$)]
  \label{prop:rat-twoadic}
  \uses{def:wseq, lem:li-mod8-2}
  \lean{QuadraticIterates.TwoAdicClass, QuadraticIterates.reflNum_emod_of_odd,
    QuadraticIterates.reflNum_emod_four_of_even,
    QuadraticIterates.not_isSquare_neg_one_zmod_reflNum,
    QuadraticIterates.not_isSquare_betaInt_of_two_le,
    QuadraticIterates.intCast_wSeq_zmod_eight_of_even, QuadraticIterates.intCast_sq_zmod_eight_of_odd,
    QuadraticIterates.intCast_pow_zmod_eight_of_even,
    QuadraticIterates.intCast_pow_two_pow_sub_one_zmod_eight_of_odd,
    Int.not_isSquare_neg_one_zmod_natAbs_of_emod}
  \leanok
  Let $N_{n_0} = w_{n_0} s^{2^{n_0-1}} + w_{n_0+1}$ (positive world: $a > 0$ or $-1 < a < 0$), resp.\
  $N_{n_0} = w^S_{n_0} s^{2^{n_0-1}} + w^S_{n_0+1}$ ($a \le -2$, $R = |r|$). If $(s \bmod 4, r \bmod 4)
  \in \{(1,1), (1,2), (3,0), (3,1)\}$ or $s$ is even and $r \equiv 3 \pmod 4$ (resp.\
  $(s, R) \bmod 4 \in \{(1,0), (1,1), (3,1), (3,2)\}$ or $s$ even and $R \equiv 3 \pmod 4$), then for
  every $n_0 \ge 2$ the integer $N_{n_0} > 0$ is divisible by $4$ or has an odd part
  $\equiv 3 \pmod 4$.
\end{proposition}
\begin{proof}
  \uses{lem:w-gamma-zmod, lem:li-period, def:betaint, lem:rat-reflection}
  \leanok
  Everything is decided modulo $8$. For odd $s$, $s^{2^{n_0}-1} \equiv s \pmod 8$, so $N \equiv
  s(\gamma_{n_0} + \gamma_{n_0+1})$ with $\gamma$ computed in $\ZZ/8\ZZ$ from $\bar a = r s^{-1}$:
  for even $\bar a$ the sequence is constant $\bar a + 1$ from $n = 2$ on
  (\lname{gammaSeq\_eq\_eval\_one}) and $N \equiv 2(r+s)$; for odd $\bar a$ it has period $2$ from
  $n = 2$ on with $\gamma_3 \in \{5, 1\}$ (\lname{gammaSeq\_eq\_ite\_even\_of\_add\_two\_eq}) and
  $N \equiv r + 2 \pmod 4$. For even $s$, $w_n \equiv r \pmod 8$ for $n \ge 3$ and $N \equiv r
  \pmod 4$. The listed classes are exactly those where the residue is as claimed (working notes,
  Proposition~5.1).
  \formalnote{In terms of $(r, s, \varepsilon)$, both worlds at once: $s$ odd with
  $r \equiv 1$ or $r + \varepsilon s \equiv 3 \pmod 4$, or $s$ even with $r \equiv 3 \pmod 4$
  (\lname{TwoAdicClass r s ε}, a \lname{Prop}, so that no signature grows). Odd $s$:
  \lname{reflNum\_emod\_of\_odd} ($N_k \equiv 6 \bmod 8$ or $N_k \equiv 3 \bmod 4$ for $k \ge 2$)
  --- in \lname{ZMod 8}, $\gamma_4 = \gamma_2$ for every $\bar a$ (a \lname{decide}), so
  \lname{gammaSeq\_eq\_ite\_even\_of\_add\_two\_eq} gives $\gamma_k + \gamma_{k+1} = \gamma_2 +
  \gamma_3$, and $N_k \equiv s(\gamma_2 + \gamma_3)$ by \lname{intCast\_wSeq} with $s^{2^k - 1} = s$;
  the inverse of $\bar s$ is $\bar s$ itself, and the residue claim is a \lname{decide} over
  $(\bar r, \bar s, \bar\varepsilon) \in (\ZZ/8\ZZ)^3$. Even $s$:
  \lname{intCast\_wSeq\_zmod\_eight\_of\_even} ($w_n \equiv r \bmod 8$ for $n \ge 3$) and
  \lname{reflNum\_emod\_four\_of\_even}. Then \lname{not\_isSquare\_neg\_one\_zmod\_reflNum}
  ($-1$ is not a square modulo $N_k$, via \lname{Int.not\_isSquare\_neg\_one\_zmod\_natAbs\_of\_emod})
  and the corollary \lname{not\_isSquare\_betaInt\_of\_two\_le}: in the classes, $\beta_n$ is not
  a square for $n \ge 2$ with $n / \rad(n) \ge 2$, by the reflection lemma with the modulus $N_k$
  itself (Mathlib's \lname{ZMod.coe\_int\_isUnit\_iff\_isCoprime} supplies the inverse of $\bar s$).}
\end{proof}

\section{Squarefree levels}

\begin{lemma}[primes of $s$]
  \label{lem:rat-ps}
  \uses{def:wseq, def:betaint}
  \lean{QuadraticIterates.intCast_wSeq_zmod_of_dvd,
    QuadraticIterates.not_isSquare_betaInt_of_squarefree_of_dvd_of_not_isSquare}
  \leanok
  For an odd prime $p \mid s$: $w_n \equiv r^{2^{n-1}-1} \pmod p$ and $\beta_n \equiv r^{E(n)}
  \pmod p$. If $\bigl(\tfrac{r}{p}\bigr) = -1$, then $\beta_n$ is a non-square modulo $p$ for every
  squarefree $n \ge 2$ (with $|r|$ in place of $r$ and $w^S$ in place of $w$ when $a \le -2$).
\end{lemma}
\begin{proof}
  \uses{lem:betaint, lem:moebiussign, lem:zmodsquare}
  \leanok
  $w_{n+1} \equiv r w_n^2 \pmod p$; $E(n)$ is odd for squarefree $n \ge 2$.
  \formalnote{\lname{intCast\_wSeq\_zmod\_of\_dvd}: $w_n \equiv r^{2^{n-1}-1}$ modulo any $M \mid s$,
  then \lname{not\_isSquare\_betaInt\_of\_squarefree\_of\_dvd\_of\_not\_isSquare}, for any modulus
  $M \mid s$ with $r$ not a square in \lname{ZMod M} (primality is not needed): $\beta_n = P/Q$
  with $P \equiv r^{F^+}$, $Q \equiv r^{F^-}$, and $F^+ + F^-$ is the odd twist exponent, by
  \lname{ZMod.isSquare\_mul\_of\_isSquare\_div} and
  \lname{odd\_sum\_two\_pow\_sub\_one\_of\_squarefree}.}
\end{proof}

\begin{lemma}[reflection at $n_0 = 1$]
  \label{lem:rat-refl1}
  \uses{lem:rat-reflection}
  \lean{QuadraticIterates.not_isSquare_betaInt_of_squarefree_of_dvd}
  \leanok
  Positive world: $\gamma_1 + \gamma_2$ has numerator $r + 2s$; if an odd prime $p \mid r + 2s$ has
  $\bigl(\tfrac{-s}{p}\bigr) = -1$, or $4 \mid r + 2s$ and $s \equiv 1 \pmod 4$, then $\beta_n$ is
  not a square for every squarefree $n \ge 2$. For $a \le -2$ the numerator is $R$, with the same
  conclusion under $\bigl(\tfrac{-s}{p}\bigr) = -1$ for an odd $p \mid R$, or $4 \mid R$ and
  $s \equiv 1 \pmod 4$.
\end{lemma}
\begin{proof}
  \uses{lem:rat-reflection}
  \leanok
  Lemma~\ref{lem:rat-reflection} with $n_0 = 1$: $\beta_n \equiv -s u^2$, and $-s$ is a non-square
  modulo $p$ iff $\bigl(\tfrac{-s}{p}\bigr) = -1$, modulo $4$ iff $s \equiv 1 \pmod 4$.
\end{proof}

\begin{lemma}[the modulus $8$ for even $s$]
  \label{lem:rat-mod8}
  \uses{def:wseq, def:betaint}
  \lean{QuadraticIterates.intCast_wSeq_zmod_eight_of_even,
    QuadraticIterates.not_isSquare_betaInt_of_squarefree_of_odd_of_even,
    QuadraticIterates.not_isSquare_betaInt_of_squarefree_of_even_of_even,
    QuadraticIterates.not_isSquare_intCast_zmod_eight_of_odd_of_emod_ne_one,
    Int.isSquare_mul_of_isSquare_div}
  \leanok
  Let $s$ be even and $r \equiv 3 \pmod 4$. Then $w_n \equiv r \pmod 8$ for $n \ge 3$, $w_2 = r + s$,
  and for squarefree $n \ge 2$, $\beta_n \equiv r^{\pm 1}$ ($n$ odd) resp.\ $(r+s)^{\pm 1}$ ($n$ even)
  modulo $8$. So $\beta_n$ is a non-square for odd $n \ge 3$, and for even $n$ if $r + s \not\equiv
  1 \pmod 8$ (with $|r|$, $|r| - s$ when $a \le -2$).
\end{lemma}
\begin{proof}
  \uses{lem:betaint, lem:moebiussign}
  \leanok
  $s^{2^n-1} \equiv 0 \pmod 8$ for $n \ge 2$ and $w_n$ is odd. Instead of the exponent of $r$ in
  $\beta_n$ one can look at the product $PQ$ of the numerator and the denominator of $\beta_n$,
  which is a square if $\beta_n$ is, and is $\prod_{d \mid n} w_d \equiv r^{\tau(n) - 1}$ (odd $n$)
  resp.\ $(r + \varepsilon s)\, r^{\tau(n) - 2}$ (even $n$) modulo $8$, with $\tau(n)$ even.
  \formalnote{\lname{intCast\_wSeq\_zmod\_eight\_of\_even}: $w_n \equiv r$ in \lname{ZMod 8} for
  $n \ge 3$ and even $s$; \lname{Int.isSquare\_mul\_of\_isSquare\_div} ($P/Q$ a rational square
  makes the integer $PQ$ a square) and \lname{card\_divisors\_eq\_two\_mul\_of\_squarefree}; the
  two cases are \lname{not\_isSquare\_betaInt\_of\_squarefree\_of\_odd\_of\_even} ($n$ odd,
  $r \equiv 3 \bmod 4$) and \lname{not\_isSquare\_betaInt\_of\_squarefree\_of\_even\_of\_even}
  ($n$ even, $r + \varepsilon s \not\equiv 1 \bmod 8$; the condition $r \equiv 3 \bmod 4$ is not
  needed for even $n$), with
  \lname{not\_isSquare\_intCast\_zmod\_eight\_of\_odd\_of\_emod\_ne\_one} (an odd integer
  $\not\equiv 1 \bmod 8$ is not a square modulo $8$).}
\end{proof}

\section{The theorems}

\begin{theorem}[Theorem R${}^+$]
  \label{thm:rat-pos}
  \uses{def:galois, def:wreath, def:wseq}
  Let $a = r/s$ with $a > 0$ or $-1 < a < 0$, $-rs$ not a square, $(r, s)$ in a class of
  Proposition~\ref{prop:rat-twoadic} (positive world), and suppose that one of (SQ${}^+$1)
  some odd $p \mid s$ has $\bigl(\tfrac{r}{p}\bigr) = -1$; (SQ${}^+$2) some odd $p \mid r + 2s$
  has $\bigl(\tfrac{-s}{p}\bigr) = -1$; (SQ${}^+$3) $4 \mid r + 2s$ and $s \equiv 1 \pmod 4$;
  (SQ${}^+$4) $s$ even and $r + s \not\equiv 1 \pmod 8$; (SQ${}^+$5) $r + s$ is not a square,
  and some odd $p \mid r + s$ has $\bigl(\tfrac{s}{p}\bigr) = -1$, or $4 \mid r + s$ and
  $s \equiv 3 \pmod 4$ holds. Then $\Om_n \cong \Ct{n}$ for all $n \ge 1$.
\end{theorem}
\begin{proof}
  \uses{lem:rat-shared, lem:rat-reflection, prop:rat-twoadic, lem:rat-ps, lem:rat-refl1, lem:rat-mod8,
    lem:rat-g2}
  All $w_n > 0$ in the positive world, so $\beta_n > 0$. For $n$ not squarefree,
  Proposition~\ref{prop:rat-twoadic} and Lemma~\ref{lem:rat-reflection}; for squarefree $n$,
  Lemmas \ref{lem:rat-ps}, \ref{lem:rat-refl1}, \ref{lem:rat-mod8}, and for (SQ${}^+$5)
  Lemma~\ref{lem:rat-g2} with $\varepsilon = 1$ ($n \ge 3$; $\beta_2 = w_2 = r + s$ is not a
  square by hypothesis); then Lemma~\ref{lem:rat-shared}.
  \formalnote{\lname{nonempty\_mulEquiv\_of\_ratClass\_of\_squarefreeCert} with the two
  \lname{Prop}-structures \lname{RatClass} and \lname{SquarefreeCert} in $r$ and $s$ (the five
  clauses as a disjunction); the sign condition $0 < r$ or $-s < r$, and $-rs$ not a square.
  Clause (SQ${}^+$5) was found during the formalization of Lemma~\ref{lem:rat-g2}: its
  mechanism works for both signs, and for $\varepsilon = 1$ it is void only for $s = 1$.
  Positivity of $w$: \lname{wSeq\_pos} for $r > 0$, and for $-s < r < 0$ the bound
  $0 < w_n \le s^{2^{n-1}-1}$.}
\end{proof}

\begin{theorem}[Theorem R${}^-$]
  \label{thm:rat-neg}
  \uses{def:galois, def:wreath, def:wseq}
  Let $a = r/s \le -2$, $R = |r|$, $-rs$ not a square, $(R, s)$ in a class of
  Proposition~\ref{prop:rat-twoadic} ($a \le -2$), and suppose that one of (SQ${}^-$1) some odd
  $p \mid s$ has $\bigl(\tfrac{R}{p}\bigr) = -1$; (SQ${}^-$2) some odd $p \mid R$ has
  $\bigl(\tfrac{-s}{p}\bigr) = -1$; (SQ${}^-$3) $4 \mid R$ and $s \equiv 1 \pmod 4$; (SQ${}^-$4)
  $R - s$ is not a square, and some odd $p \mid R - s$ has $\bigl(\tfrac{-s}{p}\bigr) = -1$ or
  $4 \mid R - s$ and $s \equiv 1 \pmod 4$ holds. Then $\Om_n \cong \Ct{n}$ for all $n \ge 1$.
\end{theorem}
\begin{proof}
  \uses{lem:rat-shared, lem:rat-reflection, prop:rat-twoadic, lem:rat-ps, lem:rat-refl1, lem:rat-g2}
  As for Theorem~\ref{thm:rat-pos} in Stoll's world ($w^S$ positive from $n = 2$ on, $\beta^S_n =
  |\beta_n|$), clause (SQ${}^-$4) by Lemma~\ref{lem:rat-g2}.
  \formalnote{Second instance of the same lemmas with $w^S$; the two worlds share every statement
  through the general \lname{gammaSeq} lemmas over \lname{ZMod M}, specialized only at the end.
  Clause (SQ${}^-$4) is Lemma~\ref{lem:rat-g2} with $\varepsilon = -1$ for $n \ge 3$, and
  $\beta^S_2 = w^S_2 = |r| - s$ is not a square by hypothesis.}
\end{proof}

\begin{lemma}[the $\gamma_2$ route, rational form of {\cite[Lemmas 3.2, 3.3]{li2020}}]
  \label{lem:rat-g2}
  \uses{def:wseq, def:betaint, lem:betaint, lem:gamma-away, lem:w-gamma-zmod, lem:li-modsq,
    lem:moebiussign}
  \lean{QuadraticIterates.not_isSquare_betaInt_of_squarefree_of_dvd_wSeq_two,
    QuadraticIterates.wSeq_two_dvd_wSeq_of_even,
    QuadraticIterates.dvd_wSeq_ediv_wSeq_two_sub_pow_of_even}
  \leanok
  In either world, let $M$ be coprime to $s$ and divide $w_2 = r + \varepsilon s$, the numerator
  of $\gamma_2$. Then $\beta_n \equiv \varepsilon s \cdot u^2 \pmod M$ with a unit $u$, for
  every squarefree $n \ge 3$; so if $\varepsilon s$ is a non-square modulo $M$, then $\beta_n$ is
  a non-square modulo $M$ for every squarefree $n \ge 3$. For $a \le -2$ ($\varepsilon = -1$,
  $M \mid |r| - s$, $-s$ a non-square) this is clause (SQ${}^-$4); in the positive world
  ($\varepsilon = 1$, $M \mid r + s$, $s$ a non-square) it is the new clause (SQ${}^+$5), void
  for $s = 1$.
\end{lemma}
\begin{proof}
  \uses{lem:li-modsq, lem:betaint, lem:gamma-away, lem:moebiussign}
  \leanok
  Over $\ZZ[1/s]$, where $M \mid \gamma_2$, Lemma~\ref{lem:li-modsq} gives $\gamma_t \equiv
  \gamma_2 \pmod{\gamma_2^2}$ for even $t$ and $\gamma_t \equiv g(0) = \varepsilon$ for odd
  $t \ge 3$. So for even $t$, $\gamma_t = \gamma_2 u_t$ with $u_t \equiv 1 \pmod M$, which in
  terms of $w$ says $w_2 \mid w_t$ in $\ZZ$ and $w_t / w_2 \equiv s^{2^{t-1} - 2} \pmod M$;
  for odd $t \ge 3$, $w_t \equiv \varepsilon s^{2^{t-1}-1}$. Write $\beta_n = P/Q$ with $P$,
  $Q$ the products of the $w_t$ over the two sign classes of the divisors $t$ of $n$. Both classes
  contain the same number $N$ of even divisors (Lemma~\ref{lem:moebiussign} on the even
  divisors), so $P = w_2^N P'$ and $Q = w_2^N Q'$ with $P'$, $Q'$ the products of the $w_2$-free
  parts, and $\beta_n = P'/Q'$. If $\beta_n$ is a square, so is the integer $P'Q'$, the product
  of the $w_2$-free parts over all divisors, which modulo $M$ is $u^{\#\{t \text{ even}\}}
  \varepsilon^{\#\{t \text{ odd}\} - 1} s^{E(n)}$ with $u = s^{-1}$; $n$ has evenly many even
  and evenly many odd divisors, and $E(n)$ is odd, so this is $\varepsilon s \cdot
  (\text{unit})^2$.
  \formalnote{\lname{not\_isSquare\_betaInt\_of\_squarefree\_of\_dvd\_wSeq\_two}, for a general
  sign $\varepsilon$, with the modulus a divisor of \lname{wSeq r s ε 2} and the hypothesis
  ``$\varepsilon s$ is not a square in \lname{ZMod M}''. The descent from $\ZZ[1/s]$
  (\lname{wSeq\_two\_dvd\_wSeq\_of\_even},
  \lname{dvd\_wSeq\_ediv\_wSeq\_two\_sub\_pow\_of\_even}) applies
  \lname{gammaSeq\_two\_sq\_dvd\_sub\_ite\_even} over \lname{Localization.Away s} and the
  descent lemma of Lemma~\ref{lem:betaint}; the $w_2$-free part is a private device as in the
  integer proof (\lname{not\_isSquare\_betaSeq\_of\_squarefree\_of\_not\_isSquare\_neg\_one}
  of Chapter~5), and the parities of the numbers of even and of odd divisors are
  \lname{Squarefree.even\_card\_filter\_divisors\_mod\_two}.}
\end{proof}

\section{Order of work}

\begin{enumerate}
\item The Galois side for $a \in \QQ$ (Theorem~\ref{thm:rat-section1}), done in one commit for
  the four files \lname{Iterates.lean} (polymorphic \lname{cSeq}, \lname{bSeq}; rational
  \lname{splittingField}, \lname{GaloisGroup}; cast plumbing), \lname{Irreducibility.lean}
  (Lemma~1.2 general, growth integral), \lname{DegreeCriterion.lean} (hypothesis swap),
  \lname{Main.lean} (Section~1 over $\QQ$, integer corollaries), together with the comparator
  challenge, after the author's approval of the change of \lname{GaloisGroup}.
\item The supporting lemmas, in the mirror files: the domain instance and the transfer of
  primes and of divisibility for \lname{Localization.Away} (Lemmas \ref{lem:gamma-away},
  \ref{lem:betaint}) in \lname{Mathlib/RingTheory/Localization/Away/Basic.lean}; the
  \lname{\_of\_dvd} primary forms and \lname{map\_moebiusFactorR} in
  \lname{Mathlib/RingTheory/MoebiusFactor.lean} (Definition~\ref{def:betaint}).
\item \lname{Rational/Sequence.lean}: Definition~\ref{def:wseq} with the identity over a general
  ring, Lemma~\ref{lem:w-coprime}, Lemma~\ref{lem:gamma-away} (with the local instances),
  Definition~\ref{def:betaint}, Lemma~\ref{lem:betaint}.
\item \lname{Rational/SquareClasses.lean}: Lemmas \ref{lem:rat-sqclass}, \ref{lem:rat-shared}
  (with the new integer lemma in \lname{Squares.lean}).
\item \lname{Rational/Reflection.lean}: Lemmas \ref{lem:w-gamma-zmod}, \ref{lem:rat-reflection}.
\item \lname{Rational/Residues.lean}: Proposition~\ref{prop:rat-twoadic}, Lemmas \ref{lem:rat-ps},
  \ref{lem:rat-refl1}, \ref{lem:rat-mod8}.
\item \lname{Rational/GammaTwo.lean}: Lemma~\ref{lem:rat-g2} (done before the Galois side, as it
  is independent of it).
\item \lname{Rational/Main.lean}: Theorem~\ref{thm:rat-pos}, then Theorem~\ref{thm:rat-neg}.
\end{enumerate}
